\documentclass[11pt]{article}

% Packages
\usepackage[utf8]{inputenc}
\usepackage[T1]{fontenc}
\usepackage{amsmath,amssymb,amsthm}
\usepackage{mathtools}
\usepackage{hyperref}
\usepackage{cleveref}
\usepackage{booktabs}
\usepackage{enumitem}
\usepackage[margin=1in]{geometry}
\usepackage{xcolor}

% Theorem environments
\newtheorem{theorem}{Theorem}[section]
\newtheorem{lemma}[theorem]{Lemma}
\newtheorem{corollary}[theorem]{Corollary}
\newtheorem{proposition}[theorem]{Proposition}
\newtheorem{definition}[theorem]{Definition}
\newtheorem{remark}[theorem]{Remark}
\newtheorem{example}[theorem]{Example}
\newtheorem{counterexample}[theorem]{Counterexample}

% Commands
\newcommand{\R}{\mathbb{R}}
\newcommand{\N}{\mathbb{N}}
\newcommand{\cl}[1]{\overline{#1}}
\newcommand{\Fix}{\operatorname{Fix}}
\newcommand{\AD}{\operatorname{AD}}
\newcommand{\id}{\operatorname{id}}
\newcommand{\diam}{\operatorname{diam}}
\newcommand{\dist}{\operatorname{dist}}

\title{Geometric Limits of Defense Design:\\
  Impossibility Theorems for Prompt Injection Defenses}

\author{
  Manish Bhatt\thanks{Corresponding author: \texttt{manish.bhatt13212@gmail.com}. This work was conducted independently and does not reflect the views of the authors' respective employers.} \\
  \and
  Sarthak Munshi \\
  \and
  Vineeth Sai Narajala \\
  \and
  Idan Habler \\
  \and
  Ammar Al-Kahfah \\
  \and
  Ken Huang \\
  \and
  Blake Gatto \\
}

\date{}

\begin{document}

\maketitle

\begin{abstract}
Prior work on LLM safety has focused on discovering adversarial inputs
or bounding perturbation robustness around individual examples. We ask
a different question: what are the \emph{geometric limits of defense
design itself}?

We prove impossibility results via two independent paths---continuous
and discrete---so that neither relies on the other.

\textbf{Path~A (Continuous).} Three results of increasing strength:
(1)~any continuous, utility-preserving defense on a connected Hausdorff
space must fix boundary points; (2)~when both defense and alignment
function are Lipschitz, the defense cannot uniformly reduce alignment
deviation far below threshold on a positive-measure band; (3)~under a
transversality condition, the defense leaves a positive-measure region
genuinely unsafe: $f(D(x)) > \tau$ persists over volume.

\textbf{Path~B (Discrete).} On finite ordered sets with no topology:
(1)~boundary crossings exist by a discrete IVT; (2)~utility-preserving
defenses must leave the safe endpoint of each crossing fixed; (3)~any
defense with bounded capacity is overwhelmed by the exponentially growing
attack surface.

The Tietze extension theorem bridges the two paths, but Path~B stands
alone for reviewers who object to continuous relaxation. The full
theory---36 files, $\sim$300 theorems, zero \texttt{sorry}, three
standard axioms---is machine-verified in Lean~4 with Mathlib and
validated against three frontier LLMs.
\end{abstract}

% ============================================================
\section{Introduction}
\label{sec:intro}

Can you build a wrapper around a language model that eliminates all
prompt injection vulnerabilities?

The defense mechanisms people have proposed---input classifiers~\cite{alon2023detecting}, output filters~\cite{inan2023llama},
constitutional rewriting~\cite{bai2022constitutional}---all share the
same structure: a function $D\colon X \to X$ that preprocesses prompts,
aiming to map unsafe inputs to safe equivalents while leaving safe
inputs alone.

Prior work has focused on discovering or mitigating individual
adversarial inputs. We ask a different question: what are the
\emph{geometric limits of the defense mapping itself}? We show that
under mild design constraints, no defense can eliminate the safety
boundary. Vulnerabilities arise not from specific inputs but from the
structure of the space and the constraints on admissible defenses.

We prove this via two independent paths, so that neither relies on the
other.

\paragraph{Path~A: Continuous impossibility (three tiers).}

\emph{Tier~1 (Boundary fixation).}
If $D$ is continuous, utility-preserving, and the space is connected,
$D$ must fix boundary points (\Cref{thm:main}).

\emph{Tier~2 ($\varepsilon$-robust constraint).}
If $D$ is $K$-Lipschitz and $f$ is $L$-Lipschitz:
\begin{equation}\label{eq:eps-bound}
  f(D(x)) \;\geq\; \tau - L(K+1)\,\delta
\end{equation}
for $x$ within distance $\delta$ of a fixed boundary point
(\Cref{thm:eps-robust}). The $\varepsilon$-band has positive measure
(\Cref{thm:band-measure}).

\emph{Tier~3 (Persistent unsafe region).}
Under transversality (boundary slope $> L(K+1)$), a positive-measure
region satisfies $f(D(x)) > \tau$---genuinely unsafe after defense
(\Cref{thm:persistent}).

\paragraph{Path~B: Discrete impossibility (no topology).}

\emph{Discrete IVT.} On $\{0,\ldots,n\}$, if $f(0) < \tau \leq f(n)$,
consecutive indices cross $\tau$ (\Cref{thm:disc-ivt}).

\emph{Discrete defense fixation.} If $D$ fixes safe vertices, the safe
side of each crossing is fixed (\Cref{thm:disc-defense}).

\emph{Capacity exhaustion.} Any defense with $< N$ distinct behaviors
misclassifies at least one of $N$ adversarial configurations
(\Cref{thm:pigeonhole}). The attack surface grows as $2^d$,
overwhelming any fixed budget (\Cref{thm:budget}).

\paragraph{Bridging the paths.}
The Tietze extension theorem (\Cref{thm:tietze}) shows that discrete
data admits continuous extensions for which Path~A holds. But Path~B
stands alone: it requires no continuous relaxation, no Tietze, no
topology.

\paragraph{Why continuity?}
Continuity captures a natural requirement: small input changes should
not produce arbitrarily large output changes. Any defense based on
learned representations or differentiable processing satisfies this.
Discontinuous defenses (blocklists, keyword filters) escape Path~A
but not Path~B.

\paragraph{Contributions.}
\begin{itemize}[nosep]
  \item \textbf{Path~A}: Three-tier continuous impossibility
    (\Cref{sec:proof,sec:eps-robust,sec:persistent}).
  \item \textbf{Path~B}: Discrete impossibility without topology
    (\Cref{sec:discrete}).
  \item Tietze bridge connecting the two (\Cref{sec:bridge}).
  \item Counterexamples proving each hypothesis necessary
    (\Cref{sec:counterexamples}).
  \item Supporting theory: basins, convergence, transferability,
    authority, gradient, stability, cost
    (\Cref{sec:landscape,sec:attacks,sec:stability,sec:cost}).
  \item Empirical validation on three LLMs (\Cref{sec:experiments}).
  \item Lean~4 artifact: 36 files, $\sim$300 theorems, zero
    \texttt{sorry} (\Cref{sec:artifact}).
\end{itemize}


% ============================================================
\section{Related Work}
\label{sec:related}

\paragraph{Adversarial robustness.}
Unlike adversarial robustness results~\cite{szegedy2013intriguing,goodfellow2014explaining,cohen2019certified},
which bound perturbation radii around individual inputs, our result is a
\emph{structural constraint on the defense mapping itself}, independent
of perturbation models, norms, or training procedures. Certified
defenses~\cite{singh2019abstract,katz2017reluplex} guarantee local
neighborhoods; we constrain \emph{global} defense behavior.

\paragraph{Prompt injection.}
Attack methods~\cite{zou2023universal,chao2024jailbreaking,mehrotra2024tree,samvelyan2024rainbow}
discover individual failures. The Manifold of Failure
framework~\cite{bhatt2025manifold} maps their topology. Our result
shows these failures are not incidental but \emph{structurally
unavoidable}.

\paragraph{Impossibility results.}
No-free-lunch theorems~\cite{wolpert1997no} limit average performance
across tasks. Robustness--accuracy tradeoffs~\cite{tsipras2018robustness}
are distributional. Our result holds \emph{pointwise} for any
continuous defense on a \emph{fixed task}, and the discrete path holds
without any continuity assumption at all.

\paragraph{Topological views.}
While our continuous proof uses elementary topology, the novelty is
identifying that \emph{standard design constraints induce a forced
fixed-point structure}. The persistent unsafe region (Tier~3) and the
discrete path go beyond topology entirely.

\paragraph{Formal verification.}
To our knowledge, this is the first formally verified impossibility
result for prompt injection defense~\cite{bagnall2019certifying,bolte2024}.


% ============================================================
\section{Setup}
\label{sec:setup}

\begin{definition}[Alignment Deviation]
\label{def:ad}
A continuous map $f\colon X \to \R$ with threshold $\tau$ partitions
$X$ into $S_\tau = \{f < \tau\}$ (safe), $U_\tau = \{f > \tau\}$
(unsafe), and $B_\tau = \{f = \tau\}$ (boundary).
\end{definition}

\begin{definition}[Defense]
\label{def:defense}
A continuous map $D\colon X \to X$ is \emph{utility-preserving} if
$D|_{S_\tau} = \id$, and \emph{complete} if $f(D(x)) < \tau$ for all~$x$.
\end{definition}


% ============================================================
\section{Path~A, Tier~1: Boundary Fixation}
\label{sec:proof}

\begin{theorem}[Boundary Fixation]
\label{thm:main}
Let $X$ be a connected Hausdorff space, $f\colon X\to\R$ continuous
with $S_\tau, U_\tau \neq \emptyset$, and $D\colon X \to X$ continuous
with $D|_{S_\tau} = \id$. Then $\exists\, z$ with $f(z) = \tau$ and
$D(z) = z$. Moreover, every $z \in \cl{S_\tau} \setminus S_\tau$
satisfies this. This set is nonempty.
\end{theorem}

\begin{proof}
(1)~$\Fix(D)$ is closed (Hausdorff $\Rightarrow$ diagonal closed).
(2)~$S_\tau \subseteq \Fix(D)$; taking closures,
$\cl{S_\tau} \subseteq \Fix(D)$.
(3)~$S_\tau$ is open and not closed (else clopen in connected space,
contradiction).
(4)~$\cl{S_\tau} \supsetneq S_\tau$; pick
$z \in \cl{S_\tau} \setminus S_\tau$: $f(z) = \tau$.
(5)~$z \in \Fix(D)$, so $D(z) = z$.
\end{proof}

\begin{corollary}
No continuous utility-preserving defense is complete.
\end{corollary}


\subsection{Each Hypothesis Is Necessary}
\label{sec:counterexamples}

\begin{counterexample}[Disconnected]
$X = \{0,1\}$ discrete: $D(1)=0$ is complete.
\end{counterexample}

\begin{counterexample}[Discontinuous]
$D(x)=x$ for $x<0.5$, $D(x)=0$ for $x \geq 0.5$: complete.
\end{counterexample}

\begin{counterexample}[Not utility-preserving]
$D(x) = x_0$ (constant): complete.
\end{counterexample}

\begin{remark}[Defense trilemma]
At most two of: continuity, utility preservation, completeness.
\end{remark}


% ============================================================
\section{Path~A, Tier~2: $\varepsilon$-Robust Constraint}
\label{sec:eps-robust}

\begin{theorem}[$\varepsilon$-Robust Constraint]
\label{thm:eps-robust}
Under \Cref{thm:main}'s hypotheses, if $f$ is $L$-Lipschitz and $D$ is
$K$-Lipschitz, the fixed point $z$ satisfies: for all $x$ (strongest
near~$z$),
$f(D(x)) \geq \tau - L(K+1)\dist(x,z)$.
\end{theorem}

\begin{theorem}[Positive-Measure Band]
\label{thm:band-measure}
The $\varepsilon$-band $\{\tau - \varepsilon \leq f \leq \tau\}$ has
positive measure (under any measure positive on nonempty open sets).
It contains $\operatorname{ball}(c, \varepsilon/(4L))$ for $c$ with
$f(c) = \tau - \varepsilon/2$.
\end{theorem}

\begin{remark}[Depth, not direction]
Tier~2 constrains how \emph{far below} $\tau$ the defense can push
near-boundary points. It does not prove they remain unsafe---the defense
may map them to $[\tau - L(K+1)\delta,\;\tau)$.
\end{remark}


% ============================================================
\section{Path~A, Tier~3: Persistent Unsafe Region}
\label{sec:persistent}

\begin{definition}[Steep region]
$\mathcal{S} = \{x : f(x) > \tau + L(K+1)\dist(x,z)\}$---points where
alignment deviation exceeds~$\tau$ by more than the defense can compensate.
\end{definition}

\begin{theorem}[Persistent Unsafe Region]
\label{thm:persistent}
If $\mathcal{S} \neq \emptyset$: (1)~$\mathcal{S}$ is open;
(2)~$\mu(\mathcal{S}) > 0$; (3)~$f(D(x)) > \tau$ for all
$x \in \mathcal{S}$.
\end{theorem}

\begin{proof}
(1)~$x \mapsto f(x) - L(K+1)\dist(x,z)$ is continuous; $\mathcal{S}$
is its superlevel set at~$\tau$.
(2)~Open $+$ nonempty $\Rightarrow$ positive measure.
(3)~$f(D(x)) \geq f(x) - L(K+1)\dist(x,z) > \tau$ for
$x \in \mathcal{S}$.
\end{proof}

\begin{remark}[Transversality]
$\mathcal{S} \neq \emptyset$ whenever the boundary has directional slope
$> L(K+1)$---the surface rises faster than the defense can pull it down.
Verified in Lean: if $f$ has directional derivative $c > L(K+1)$ along
unit vector~$v$ at~$z$, then $z + tv \in \mathcal{S}$ for small $t > 0$.
\end{remark}


% ============================================================
\section{Path~B: Discrete Impossibility}
\label{sec:discrete}

Path~B requires no topology, no continuity, no Tietze extension.

\begin{theorem}[Discrete IVT]
\label{thm:disc-ivt}
Let $f\colon \{0,\ldots,n\!+\!1\} \to \R$ with $f(0) < \tau$ and
$f(n\!+\!1) \geq \tau$. Then $\exists\, i$ with $f(i) < \tau$ and
$f(i\!+\!1) \geq \tau$.
\end{theorem}

\begin{proof}
By induction: if $f(0) < \tau$ and no consecutive crossing exists,
then $f(k) < \tau$ for all $k$, contradicting $f(n+1) \geq \tau$.
\end{proof}

\begin{theorem}[Discrete Defense Fixation]
\label{thm:disc-defense}
If $D$ fixes safe vertices ($f(v) < \tau \Rightarrow D(v) = v$),
the crossing from \Cref{thm:disc-ivt} has its safe endpoint fixed.
\end{theorem}

\begin{theorem}[Capacity Exhaustion]
\label{thm:pigeonhole}
If there are $m$ adversarial configurations and the defense realizes
$< m$ distinct behaviors, at least two configurations receive the same
treatment (pigeonhole).
\end{theorem}

\begin{theorem}[Exponential Attack Surface]
\label{thm:budget}
$2^d \geq d + 1$ for all $d$. For any fixed budget $B$, there exists
$d_0$ with $2^d > B$ for all $d \geq d_0$. The attack surface grows
exponentially; no fixed-capacity defense suffices.
\end{theorem}

\begin{remark}[Discrete trilemma]
On a finite ordered set: if $f$ crosses $\tau$ and $D$ fixes safe
inputs, $D$ must leave the safe endpoint of the crossing fixed. This
is the discrete boundary fixation, proved by induction rather than
topology.
\end{remark}


% ============================================================
\section{Bridge: Tietze Extension}
\label{sec:bridge}

\begin{theorem}[Continuous Relaxation]
\label{thm:tietze}
Let $X$ be connected, normal, Hausdorff with finite observations
$g(p) < \tau < g(q)$. Then $\exists$ continuous $f$ agreeing with all
observations for which Path~A holds. Lipschitz extensions enable
Tiers~2--3.
\end{theorem}

\begin{remark}
Path~B does not need \Cref{thm:tietze}. The bridge connects the two
paths but neither depends on the other for its core claims.
\end{remark}


% ============================================================
\section{Vulnerability Landscape}
\label{sec:landscape}

\begin{theorem}[Basin Structure]
\label{thm:basin}
$U_\tau$ is open; positive measure under any measure positive on
nonempty open sets.
\end{theorem}

\begin{theorem}[Fragment Size]
\label{thm:fragment}
$L$-Lipschitz $f$ with $f(p) > \tau$: connected component of $U_\tau$
at $p$ has diameter $\geq 2(f(p)-\tau)/L$.
\end{theorem}


% ============================================================
\section{Attack Properties}
\label{sec:attacks}

\begin{theorem}[Perturbation Robustness]
\label{thm:lipschitz}
$\operatorname{ball}(p, (f(p)-\tau)/L) \subseteq U_\tau$; radius
monotone in $f(p)$.
\end{theorem}

\begin{theorem}[Convergence]
\label{thm:convergence}
Monotone-improvement operators on $[0,1]$ converge in
$\leq \lfloor 1/\delta \rfloor$ steps.
\end{theorem}

\begin{theorem}[Transferability]
\label{thm:transfer}
$\|f-g\|_\infty \leq \delta \Rightarrow \{f > \tau+\delta\} \subseteq
\{g > \tau\}$. Zero-cost transfer.
\end{theorem}

\begin{theorem}[Authority Monotonicity]
\label{thm:authority}
$f(y,\cdot)$ monotone in authority $\Rightarrow$ vulnerability
upward-closed; critical threshold exists by IVT when $f(y,\cdot)$
is continuous and crosses~$\tau$. If also monotone in indirection,
threshold curve is non-increasing.
\end{theorem}

\begin{theorem}[Gradient Ascent]
\label{thm:gradient}
Nonzero Fr\'{e}chet derivative at $x$ $\Rightarrow$ $\exists\, v,
\varepsilon > 0$: $f(x + \varepsilon v) > f(x)$.
\end{theorem}


% ============================================================
\section{Stability Under Fine-Tuning}
\label{sec:stability}

\begin{theorem}[Interior Stability]
\label{thm:interior-stable}
$\|f-g\|_\infty \leq \varepsilon$: deep basin points and deep safe
points preserve classification. Only $|f(x)-\tau| \leq \varepsilon$
is uncertain.
\end{theorem}

\begin{theorem}[Crossing Preservation]
\label{thm:crossing}
Transversal crossings persist under $\varepsilon$-perturbation.
\end{theorem}

\begin{theorem}[Patching Is Nonlocal]
\label{thm:nonlocal}
Eliminating a vulnerability at one point necessarily changes values
elsewhere. (Explicit Lean construction.)
\end{theorem}


% ============================================================
\section{Cost Asymmetry}
\label{sec:cost}

\begin{theorem}[Grid-Based Cost Asymmetry]
\label{thm:cost}
Grid defense: $N^d$ cost (exponential). Attack: $\leq 1/\delta$
(dimension-free). Ratio $\delta N^d \to \infty$.
\end{theorem}


% ============================================================
\section{Experimental Validation}
\label{sec:experiments}

Validated against the Manifold of Failure
framework~\cite{bhatt2025manifold} on three LLMs over a 2D behavioral
space.

\begin{table}[t]
\centering
\small
\caption{Falsifiable predictions confirmed by empirical data.}
\label{tab:correspondence}
\begin{tabular}{@{}p{2.8cm}p{4.5cm}p{4.5cm}@{}}
\toprule
\textbf{Theorem} & \textbf{Predicts} & \textbf{Confirmed by} \\
\midrule
Basin (\ref{thm:basin})
  & Positive-measure basins
  & Extended heatmap regions \\[2pt]
Fragment (\ref{thm:fragment})
  & Smooth $\to$ large; rough $\to$ mosaic
  & Llama: mesa; GPT-OSS: mosaic \\[2pt]
Convergence (\ref{thm:convergence})
  & Monotone attack convergence
  & Curves plateau \\[2pt]
Transfer (\ref{thm:transfer})
  & Similar surfaces $\to$ shared basins
  & $.93 \to .73 \to .47$ spectrum \\[2pt]
Authority (\ref{thm:authority})
  & Horizontal banding
  & Bands at $a_2 \approx .3, .7$ \\[2pt]
Persistent (\ref{thm:persistent})
  & Steep boundary $\to$ unsafe volume
  & Llama's $.93$ plateau persists \\[2pt]
Stability (\ref{thm:interior-stable})
  & Deep points survive fine-tuning
  & Vulnerability persists across variants \\[2pt]
Cost (\ref{thm:cost})
  & $d\!=\!2$ tractable; high-$d$ not
  & 15K queries fill 63\% at $d\!=\!2$ \\
\bottomrule
\end{tabular}
\end{table}

\paragraph{Llama-3-8B} (AD $0.93$, basin rate $93.9\%$): Tier~3
territory---boundary slope far exceeds $L(K+1)$ for reasonable~$K$.

\paragraph{GPT-OSS-20B} (AD $0.73$, basin rate $64.3\%$): rugged
landscape, many transversal crossings, each generating a persistent
unsafe pocket.

\paragraph{GPT-5-Mini} (peak AD $0.50$, basin rate $0\%$): $U_{0.5} =
\emptyset$, so no tier applies---correctly predicting no impossibility.
Strategy: make the boundary harmless.


% ============================================================
\section{Engineering Prescription}
\label{sec:engineering}

\textbf{Make the boundary shallow} (raise $\tau$; GPT-5-Mini).
\textbf{Reduce} $L$ (smoother surface, predictable landscape).
\textbf{Reduce dimension} (constrain prompt interface).
\textbf{Monitor the boundary} (crossings persist under fine-tuning).


% ============================================================
\section{The Lean Artifact}
\label{sec:artifact}

Verified in Lean~4.28.0 + Mathlib~v4.28.0. 36 files:
\begin{itemize}[nosep]
  \item 10 core (MoF\_01--10), 10 cost (Cost\_01--10), 10 advanced
    (Adv\_01--10)
  \item MoF\_ContinuousRelaxation (Tietze bridge)
  \item MoF\_11\_EpsilonRobust ($\varepsilon$-robust + persistent)
  \item MoF\_12\_Discrete (discrete path)
  \item 3 capstone (MasterTheorem, Euclidean, verification)
\end{itemize}

Three axioms: \texttt{propext}, \texttt{Classical.choice},
\texttt{Quot.sound}. Zero additional axioms.


% ============================================================
\section{Limitations}
\label{sec:limitations}

\paragraph{Tier~1} is at the boundary ($f(z) = \tau$ exactly).
\paragraph{Tier~2} constrains depth, not direction.
\paragraph{Tier~3} requires transversality ($c > L(K+1)$).
\paragraph{Path~B} requires ordered structure; general graphs need
further work.
\paragraph{Tietze} applies to interpolants, not the discrete system
itself.
\paragraph{Cost asymmetry} assumes grid enumeration.
\paragraph{Single-turn.} Manifold may shift across turns.


% ============================================================
\section{Conclusion}
\label{sec:conclusion}

We have established defense impossibility via two independent paths.

\textbf{Path~A} (continuous): boundary fixation $\to$
$\varepsilon$-robust band $\to$ persistent unsafe volume under
transversality. Each tier strictly strengthens the last.

\textbf{Path~B} (discrete): discrete IVT $\to$ defense fixation $\to$
capacity exhaustion. No topology, no Tietze, no continuity.

The key insight is not about attacks or training but about the
\emph{geometry of the defense itself}: design constraints create
anchors that the defense cannot escape, whether the space is continuous
or discrete.

The practical prescription: make the boundary shallow, smooth,
and low-dimensional. GPT-5-Mini's ceiling at $\AD = 0.50$ shows this
working.

36 files, $\sim$300 theorems, zero \texttt{sorry}, three standard
axioms.


% ============================================================
\bibliographystyle{plain}
\begin{thebibliography}{99}

\bibitem{alon2023detecting}
G.~Alon and M.~Kamfonas.
Detecting language model attacks with perplexity.
\emph{arXiv:2308.14132}, 2023.

\bibitem{bagnall2019certifying}
A.~Bagnall and G.~Stewart.
Certifying the true error: ML in {Coq} with verified generalization.
\emph{AAAI}, 2019.

\bibitem{bai2022constitutional}
Y.~Bai et~al.
Constitutional {AI}: Harmlessness from {AI} feedback.
\emph{arXiv:2212.08073}, 2022.

\bibitem{bhatt2025manifold}
M.~Bhatt et~al.
Manifold of failure: Behavioral attraction basins in language models.
\emph{arXiv:2602.22291v2}, 2026.

\bibitem{bolte2024}
J.~Bolte et~al.
Formal verification of optimization algorithms.
\emph{Math.\ Programming}, 2024.

\bibitem{chao2024jailbreaking}
P.~Chao et~al.
Jailbreaking black box LLMs in twenty queries.
\emph{arXiv:2310.08419}, 2024.

\bibitem{cohen2019certified}
J.~Cohen, E.~Rosenfeld, J.~Z. Kolter.
Certified adversarial robustness via randomized smoothing.
\emph{ICML}, 2019.

\bibitem{goodfellow2014explaining}
I.~J. Goodfellow, J.~Shlens, C.~Szegedy.
Explaining and harnessing adversarial examples.
\emph{arXiv:1412.6572}, 2014.

\bibitem{inan2023llama}
H.~Inan et~al.
{Llama Guard}: {LLM}-based input-output safeguard.
\emph{arXiv:2312.06674}, 2023.

\bibitem{katz2017reluplex}
G.~Katz et~al.
Reluplex: Efficient {SMT} for verifying deep neural networks.
\emph{CAV}, 2017.

\bibitem{mehrotra2024tree}
A.~Mehrotra et~al.
Tree of attacks: Jailbreaking black-box {LLMs} automatically.
\emph{arXiv:2312.02119}, 2024.

\bibitem{samvelyan2024rainbow}
M.~Samvelyan et~al.
Rainbow teaming: Open-ended generation of diverse adversarial prompts.
\emph{NeurIPS}, 37, 2024.

\bibitem{singh2019abstract}
G.~Singh et~al.
An abstract domain for certifying neural networks.
\emph{POPL}, 2019.

\bibitem{szegedy2013intriguing}
C.~Szegedy et~al.
Intriguing properties of neural networks.
\emph{arXiv:1312.6199}, 2013.

\bibitem{tsipras2018robustness}
D.~Tsipras et~al.
Robustness may be at odds with accuracy.
\emph{arXiv:1805.12152}, 2018.

\bibitem{wolpert1997no}
D.~H. Wolpert and W.~G. Macready.
No free lunch theorems for optimization.
\emph{IEEE Trans.\ Evol.\ Comput.}, 1(1):67--82, 1997.

\bibitem{zou2023universal}
A.~Zou et~al.
Universal and transferable adversarial attacks on aligned language models.
\emph{arXiv:2307.15043}, 2023.

\end{thebibliography}

\end{document}
